\scanpage{92}
\deleted{\unsure{Le|Je}} Montrons que a) implique c). Soit en
effet \deleted{\unsure{donc}} $E$ un vectoriel \deleted{\unsure{avec $T$|sur $Q$}}\inserted{``de degré $i$''} avec $T$ opérant.
\deleted{\unsure{Alors les opé|Alors les opérations de} \unsure{Donc} $T_{\mathbb R}$ opère sur $E_{\mathbb R}$, et comme
\unsure{c'est là la chose qui se passe}}
Alors $T_{\mathbb R}$ opère sur $E_{\mathbb R}$. Soit $U$ le plus
grand sous-tore compact de $T_{\mathbb R}$. Il résulte
aisément de l'hyp. a) que $T_{\mathbb R}^{\unsure{\circ|0}}=U_{\deleted{\unsure{R}}\mathbb R}\cdot S_{\mathbb R}$,
où $S=\operatorname{Im}j$, \deleted{\unsure{et}}\inserted{et $U\subset\operatorname{Ker}\varepsilon$.} \deleted{Donc} $f$ est de la forme
$u\,j(x)$, ou du moins il existe un entier \deleted{\unsure{$\nu$ tel que}}\inserted{$\nu>0$ \unsure{tq}}
que
\[
f^\nu=u\cdot j(x)\qquad u\in U(\mathbb R),\quad x\in\mathbb R^*
\]
Or $U$ étant compact, les valeurs propres de $\forall$
$u\in U(\mathbb R)$ sont de val absolue $1$, donc celles
de $f^\nu$ sont de val. absolues égales à celles
de $j(x)$, or $j(x)$ est la \deleted{\unsure{signature}}\unsure{l'hom.} \deleted{\unsure{de}} par
$x^i$ \unsure{,} \inserted{donc les val. propres de $f$ sont de val. abs. $x^{i/\nu}$.}
D'autre part $\deleted{\varepsilon(u)=0}\ \varepsilon\mid U=0$ donc
\[
\varepsilon(\unsure{f|f^\nu})^{\deleted{\unsure{\nu}}\nu}=\varepsilon(u)\,\varepsilon j(x)=\varepsilon j(x)=x^\alpha,\qquad
\text{donc}\quad x^{i/\nu}=\varepsilon(f)^{i/\alpha}=q^{i/\alpha},
\]
d'où la conclusion.

\margin{?}
\deleted{[NB. Evidemment b) et c) n'impliquent pas a), \unsure{pour aucune|pour mieux},
car il y a des sous-groupes libres de $Q^*$ de
rang $\geq 2$ \unsure{dans|donc}]}

D'autre part, on voit que a équivaut au
fait que $\mathcal M$ est (avec ses deux structures suppl.
provenant de $j,\varepsilon$) munissable d'une structure
``positive'' du style motivique. Ceci, d'autre part,
\unsure{via|par} un argument connu, implique que
% archive sheet number: 91

